%===========================================================================
% PART 2: CFA LEVEL II STYLE - FLAT STRUCTURE (NO NESTING ERRORS)
%===========================================================================

\documentclass[aspectratio=169,10pt]{beamer}
\usetheme{Madrid}
\usecolortheme{default}
\usepackage{amsmath, amssymb, amsfonts}
\usepackage{booktabs}
\usepackage{array}
\usepackage{xcolor}
\usepackage{multicol}
\usepackage{tikz}

% Colors
\definecolor{navyblue}{RGB}{0,53,102}
\definecolor{crimson}{RGB}{165,42,42}
\definecolor{darkgreen}{RGB}{0,100,0}

% Font sizes
\renewcommand{\small}{\fontsize{6pt}{7.5pt}\selectfont}
\renewcommand{\normalsize}{\fontsize{7pt}{9pt}\selectfont}
\renewcommand{\large}{\fontsize{8pt}{10pt}\selectfont}
\renewcommand{\Large}{\fontsize{9pt}{11pt}\selectfont}

% Custom commands
\newcommand{\question}[1]{{\large\bfseries\color{crimson}#1}}
\newcommand{\subquestion}[1]{{\bfseries\color{darkgreen}#1}}
\newcommand{\highlight}[1]{\textcolor{crimson}{\textbf{#1}}}
\newcommand{\boxedanswer}[1]{\fcolorbox{crimson}{white}{\parbox{0.9\linewidth}{\small #1}}}

% Title
\title[Synoptic Exam]{Advanced Financial Theory}
\subtitle{CFA Level II Style Questions \& Solutions}
\author{Shivgan Joshi}
\date{July 8, 2026}

\begin{document}
	
	%===========================================================================
	% TITLE
	%===========================================================================
	\begin{frame}
		\titlepage
	\end{frame}
	
	%===========================================================================
	% SECTION C TITLE
	%===========================================================================
	\begin{frame}{Section C: Capital Structure}
		\begin{center}
			\Large\textcolor{navyblue}{Section C}
			\vspace{0.2cm}
			\large Capital Structure and Strategic Finance
		\end{center}
	\end{frame}
	
	%===========================================================================
	% Q7 - PROBLEM
	%===========================================================================
	\begin{frame}{Q7: Capital Structure Conundrum}
		\question{Context}
		\small
		Aethelred Industries: Equity \$3.5B, Debt \$1.5B, R$_E$=14\%, R$_D$=6\%, T$_c$=21\%
		\vspace{0.1cm}
		
		\question{Requirements}
		\begin{enumerate}
			\item Calculate R$_A$ using MM Prop II with taxes.
			\item Calculate WACC before and after tax.
			\item Calculate tax shield value.
			\item Explain trade-off and pecking order theories.
			\item Recommend optimal capital structure.
		\end{enumerate}
	\end{frame}
	
	%===========================================================================
	% Q7 - SOLUTION A
	%===========================================================================
	\begin{frame}{Q7 Solution: MM Calculations}
		\question{Part (a): Unlevered Cost of Equity}
		\small
		\[
		R_E = R_A + \frac{D}{E}(R_A - R_D)(1 - T_c)
		\]
		\[
		0.14 = R_A + 0.4286(R_A - 0.06)(0.79)
		\]
		\[
		0.14 = R_A + 0.3386R_A - 0.0203
		\]
		\[
		0.1603 = 1.3386R_A
		\]
		\[
		\boxed{R_A = 11.97\%}
		\]
		
		\vspace{0.1cm}
		\small
		The unlevered cost of equity (11.97\%) is lower than levered cost (14\%) because leverage increases equity risk.
	\end{frame}
	
	%===========================================================================
	% Q7 - SOLUTION B
	%===========================================================================
	\begin{frame}{Q7 Solution: WACC}
		\question{Part (b): WACC Calculations}
		\small
		V = 3.5 + 1.5 = \$5.0B
		
		\vspace{0.1cm}
		\textbf{WACC Before Tax:}
		\[
		WACC_{BT} = \frac{3.5}{5.0}(0.14) + \frac{1.5}{5.0}(0.06) = 0.098 + 0.018 = \boxed{11.60\%}
		\]
		
		\vspace{0.1cm}
		\textbf{WACC After Tax:}
		\[
		WACC_{AT} = \frac{3.5}{5.0}(0.14) + \frac{1.5}{5.0}(0.06)(0.79) = 0.098 + 0.0142 = \boxed{11.22\%}
		\]
		
		\vspace{0.1cm}
		\small
		After-tax WACC is lower because interest is tax-deductible.
	\end{frame}
	
	%===========================================================================
	% Q7 - SOLUTION C
	%===========================================================================
	\begin{frame}{Q7 Solution: Tax Shield}
		\question{Part (c): Tax Shield Value}
		\small
		V$_L$ = E + D = 3.5 + 1.5 = \$5.0B
		
		\vspace{0.1cm}
		\[
		V_U = V_L - T_c \times D = 5.0 - 0.21(1.5) = \boxed{\$4.685B}
		\]
		
		\vspace{0.1cm}
		\[
		PV(\text{Tax Shield}) = T_c \times D = 0.21 \times 1.5 = \boxed{\$315M}
		\]
		
		\vspace{0.1cm}
		\small
		MM Proposition I with Taxes: V$_L$ = V$_U$ + T$_c$$\times$D
		\[
		5.0 = 4.685 + 0.315
		\]
		The tax shield adds \$315M to firm value.
	\end{frame}
	
	%===========================================================================
	% Q7 - SOLUTION D
	%===========================================================================
	\begin{frame}{Q7 Solution: Trade-Off Theory}
		\question{Part (d): Trade-Off Theory}
		\small
		\[
		V_L = V_U + PV(\text{Tax Shield}) - PV(\text{Distress Costs}) - PV(\text{Agency Costs})
		\]
		
		\vspace{0.1cm}
		\textbf{Benefits of Debt:}
		\begin{itemize}
			\item Tax shield (interest is tax-deductible)
			\item Discipline on management
			\item Signals confidence
		\end{itemize}
		
		\vspace{0.1cm}
		\textbf{Costs of Debt:}
		\begin{itemize}
			\item Financial distress (bankruptcy risk)
			\item Agency costs (debt-equity conflicts)
			\item Loss of flexibility (covenants)
		\end{itemize}
	\end{frame}
	
	%===========================================================================
	% Q7 - SOLUTION E
	%===========================================================================
	\begin{frame}{Q7 Solution: Pecking Order}
		\question{Part (e): Pecking Order Theory}
		\small
		Myers \& Majluf (1984)
		
		\vspace{0.1cm}
		\textbf{Financing Hierarchy:}
		\begin{enumerate}
			\item Internal financing (retained earnings)
			\item Debt (less information-sensitive)
			\item Equity (adverse selection)
		\end{enumerate}
		
		\vspace{0.1cm}
		\textbf{Key Assumption:} Information asymmetry (managers know more).
		
		\vspace{0.1cm}
		\textbf{Implications:}
		\begin{itemize}
			\item No target debt ratio
			\item Financing driven by needs
			\item Profitable firms use less debt
			\item Equity issuance signals overvaluation
		\end{itemize}
	\end{frame}
	
	%===========================================================================
	% Q7 - SOLUTION F
	%===========================================================================
	\begin{frame}{Q7 Solution: Optimal Structure}
		\question{Part (f): Strategic Recommendation}
		\small
		\begin{table}[h]
			\centering
			\begin{tabular}{|c|c|c|c|}
				\hline
				\textbf{D/E} & \textbf{R$_D$} & \textbf{R$_E$} & \textbf{WACC} \\
				\hline
				0.43 & 6.0\% & 14.00\% & 11.22\% \\
				0.50 & 6.2\% & 14.35\% & 11.03\% \\
				0.60 & 6.5\% & 14.79\% & \textbf{10.85\%} \\
				0.70 & 7.0\% & 15.23\% & 10.96\% \\
				0.80 & 7.5\% & 15.67\% & 11.12\% \\
				\hline
			\end{tabular}
		\end{table}
		
		\vspace{0.1cm}
		\textbf{Optimal D/E = 0.60} (minimizes WACC)
		
		\vspace{0.1cm}
		\textbf{Value Creation:} Increases firm value by \$1.19B
	\end{frame}
	
	%===========================================================================
	% Q8 - PROBLEM
	%===========================================================================
	\begin{frame}{Q8: Dividend Policy Decision}
		\question{Context}
		\small
		Shares=100M, Price=\$50, Earnings=\$300M, Dividend=\$2.00, Repurchase=\$500M
		
		\vspace{0.1cm}
		\question{Requirements}
		\begin{enumerate}
			\item Compare dividend theories.
			\item Calculate repurchase impact on EPS, price, value.
			\item Analyze behavioral considerations.
			\item Recommend optimal policy.
		\end{enumerate}
	\end{frame}
	
	%===========================================================================
	% Q8 - SOLUTION A
	%===========================================================================
	\begin{frame}{Q8 Solution: Dividend Theories}
		\question{Part (a): Dividend Theories}
		\small
		
		\textbf{1. Dividend Irrelevance (MM)}
		Dividend policy does not affect firm value in perfect markets.
		
		\vspace{0.1cm}
		\textbf{2. Bird-in-the-Hand}
		Dividends are less risky than capital gains; investors prefer dividends.
		\[
		P_0 = \frac{D_1}{R_E - g}
		\]
		
		\vspace{0.1cm}
		\textbf{3. Tax Preference}
		Investors prefer capital gains (taxed at 15\% vs. dividends at 20\%).
		Lower dividends $\rightarrow$ higher after-tax returns.
	\end{frame}
	
	%===========================================================================
	% Q8 - SOLUTION B
	%===========================================================================
	\begin{frame}{Q8 Solution: Repurchase Analysis}
		\question{Part (b): Share Repurchase Analysis}
		\small
		\[
		\text{EPS} = \frac{300,000,000}{100,000,000} = \boxed{\$3.00}
		\]
		\[
		\text{P/E} = \frac{50.00}{3.00} = \boxed{16.67x}
		\]
		\[
		\text{Shares Repurchased} = \frac{500,000,000}{50.00} = \boxed{10,000,000}
		\]
		\[
		\text{New EPS} = \frac{300,000,000}{90,000,000} = \boxed{\$3.33}
		\]
		\[
		\text{New Price} = 3.33 \times 16.67 = \boxed{\$55.50}
		\]
		\textbf{Value Created:} \$4,995M - \$5,000M = \boxed{-\$5M}
	\end{frame}
	
	%===========================================================================
	% Q8 - SOLUTION C
	%===========================================================================
	\begin{frame}{Q8 Solution: Behavioral Considerations}
		\question{Part (c): Behavioral Considerations}
		\small
		
		\textbf{Mental Accounting}
		Investors view dividends as "income" and capital gains as "savings." Dividends are spent more readily.
		
		\vspace{0.1cm}
		\textbf{Self-Control}
		Dividends provide pre-commitment; retirees prefer dividends.
		
		\vspace{0.1cm}
		\textbf{Signaling}
		Dividend increases signal confidence (+2-3\% reaction).
		Dividend decreases signal pessimism (-5-8\% reaction).
	\end{frame}
	
	%===========================================================================
	% Q8 - RECOMMENDATION
	%===========================================================================
	\begin{frame}{Q8: Comprehensive Recommendation}
		\question{Recommendation}
		\small
		
		\textbf{1. Optimal Payout:}
		Reduce from 66.7\% to 50\% (\$1.50/share)
		
		\vspace{0.1cm}
		\textbf{2. Trade-off:}
		\begin{table}[h]
			\centering
			\begin{tabular}{|c|c|c|}
				\hline
				\textbf{Criteria} & \textbf{Dividend} & \textbf{Repurchase} \\
				\hline
				Tax Efficiency & Lower & Higher \\
				Flexibility & Lower & Higher \\
				Signaling & Strong & Moderate \\
				\hline
			\end{tabular}
		\end{table}
		
		\vspace{0.1cm}
		\textbf{3. Policy:}
		Increase dividend to \$2.20/share + \$200M buyback
	\end{frame}
	
	%===========================================================================
	% Q9 - PROBLEM
	%===========================================================================
	\begin{frame}{Q9: Merger and Acquisition}
		\question{Context}
		\small
		Aethelred (MC \$15B, P/E 12.5x) acquiring Mercian (MC \$4.5B, P/E 22.5x)
		
		\vspace{0.1cm}
		FCF: Y1=150, Y2=180, Y3=216, Y4=259, Y5=311, TV=8,400 (WACC=10\%)
		
		\vspace{0.1cm}
		\question{Requirements}
		\begin{enumerate}
			\item Calculate DCF value.
			\item Analyze premium scenarios and synergies.
			\item Evaluate behavioral biases.
			\item Provide recommendation.
		\end{enumerate}
	\end{frame}
	
	%===========================================================================
	% Q9 - SOLUTION DCF
	%===========================================================================
	\begin{frame}{Q9 Solution: DCF Valuation}
		\question{Part (a): DCF Valuation}
		\small
		\[
		PV(FCF) = \frac{150}{1.10} + \frac{180}{1.10^2} + \frac{216}{1.10^3} + \frac{259}{1.10^4} + \frac{311}{1.10^5}
		\]
		\[
		= 136.36 + 148.76 + 162.28 + 176.34 + 193.06 = \boxed{816.80M}
		\]
		\[
		PV(TV) = \frac{8,400}{1.10^5} = \frac{8,400}{1.6105} = \boxed{5,215.77M}
		\]
		\[
		\text{Enterprise Value} = 816.80 + 5,215.77 = \boxed{\$6,032.57M}
		\]
		\[
		\text{Equity Value} = 6,032.57 - 1,000 = \boxed{\$5,032.57M}
		\]
	\end{frame}
	
	%===========================================================================
	% Q9 - SOLUTION PREMIUM
	%===========================================================================
	\begin{frame}{Q9 Solution: Premium Analysis}
		\question{Part (b): Premium Scenarios}
		\small
		\begin{table}[h]
			\centering
			\begin{tabular}{|c|c|c|c|}
				\hline
				\textbf{Premium} & \textbf{Price} & \textbf{Value (\$M)} & \textbf{Premium} \\
				\hline
				20\% & \$54.00 & 5,400 & 900 \\
				30\% & \$58.50 & 5,850 & 1,350 \\
				40\% & \$63.00 & 6,300 & 1,800 \\
				\hline
			\end{tabular}
		\end{table}
		
		\vspace{0.1cm}
		DCF value = \$5,033M (11.8\% premium).
		Premium > 30\% destroys value.
	\end{frame}
	
	%===========================================================================
	% Q9 - SOLUTION SYNERGIES
	%===========================================================================
	\begin{frame}{Q9 Solution: Synergies}
		\question{Part (c): Synergy Analysis}
		\small
		\textbf{Revenue Synergies:} \$150M/year
		\textbf{Cost Synergies:} \$150M/year
		\textbf{Total:} \$300M/year
		
		\vspace{0.1cm}
		\[
		PV(\text{Synergies}) = \frac{300}{0.10} - 200 \text{ (integration)} = \boxed{\$2,800M}
		\]
		
		\vspace{0.1cm}
		\textbf{Total Value:} \$5,033 + \$2,800 = \$7,833M
	\end{frame}
	
	%===========================================================================
	% Q9 - SOLUTION BEHAVIORAL
	%===========================================================================
	\begin{frame}{Q9 Solution: Behavioral Factors}
		\question{Part (d): Behavioral Factors}
		\small
		
		\textbf{CEO Hubris}
		Overconfidence, empire-building
		
		\vspace{0.1cm}
		\textbf{Winner's Curse}
		Overpaying in bidding wars
		
		\vspace{0.1cm}
		\textbf{Agency Problems}
		Personal gain, firm size compensation
	\end{frame}
	
	%===========================================================================
	% Q9 - RECOMMENDATION
	%===========================================================================
	\begin{frame}{Q9: Comprehensive Recommendation}
		\question{Recommendation}
		\small
		
		\textbf{1. Proceed?} Yes, but cautiously (strong strategic fit)
		
		\vspace{0.1cm}
		\textbf{2. Max Premium:} 30\% (\$58.50/share)
		
		\vspace{0.1cm}
		\textbf{3. Financing:} 50\% cash, 50\% stock
		
		\vspace{0.1cm}
		\textbf{4. Integration:}
		\begin{itemize}
			\item Dedicated integration team
			\item 100-day review
			\item Clear KPIs
		\end{itemize}
	\end{frame}
	
	%===========================================================================
	% SECTION D TITLE
	%===========================================================================
	\begin{frame}{Section D: Corporate Governance}
		\begin{center}
			\Large\textcolor{navyblue}{Section D}
			\vspace{0.2cm}
			\large Corporate Governance and Value Creation
		\end{center}
	\end{frame}
	
	%===========================================================================
	% Q10 - PROBLEM
	%===========================================================================
	\begin{frame}{Q10: Enterprise Risk Management}
		\question{Context}
		\small
		Aethelred implementing ERM framework after market volatility.
		
		\vspace{0.1cm}
		\question{Requirements}
		\begin{enumerate}
			\item Explain ERM principles.
			\item Design ERM framework.
			\item Calculate VaR and CVaR.
			\item Discuss behavioral biases.
			\item Explain ERM as counterweight.
		\end{enumerate}
	\end{frame}
	
	%===========================================================================
	% Q10 - SOLUTION PRINCIPLES
	%===========================================================================
	\begin{frame}{Q10 Solution: ERM Principles}
		\question{Part (a): ERM Principles}
		\small
		
		\textbf{1. Integrated Risk View}
		Interconnected risks viewed holistically
		
		\vspace{0.1cm}
		\textbf{2. Portfolio Approach}
		Diversification reduces overall risk
		
		\vspace{0.1cm}
		\textbf{3. Strategic Alignment}
		Risk management supports strategy
		
		\vspace{0.1cm}
		\textbf{4. Continuous Monitoring}
		Dynamic and forward-looking
	\end{frame}
	
	%===========================================================================
	% Q10 - SOLUTION FRAMEWORK
	%===========================================================================
	\begin{frame}{Q10 Solution: ERM Framework}
		\question{Part (b): ERM Framework}
		\small
		
		\textbf{1. Risk Identification}
		Strategic, Financial, Operational, Compliance
		
		\vspace{0.1cm}
		\textbf{2. Risk Assessment}
		Quantitative (P$\times$I), Qualitative (heat maps), Scenario analysis
		
		\vspace{0.1cm}
		\textbf{3. Mitigation}
		Avoid, Reduce, Transfer, Accept
		
		\vspace{0.1cm}
		\textbf{4. Monitoring}
		KRIs, Board committee, Regular reporting
	\end{frame}
	
	%===========================================================================
	% Q10 - SOLUTION RISK QUANT
	%===========================================================================
	\begin{frame}{Q10 Solution: Risk Quantification}
		\question{Part (c): Risk Quantification}
		\small
		\begin{table}[h]
			\centering
			\begin{tabular}{|c|c|c|c|}
				\hline
				\textbf{Risk} & \textbf{Prob} & \textbf{Loss (\$M)} & \textbf{Expected Loss} \\
				\hline
				Market & 25\% & 200 & 50 \\
				Credit & 15\% & 150 & 22.5 \\
				Operational & 20\% & 100 & 20 \\
				Regulatory & 10\% & 75 & 7.5 \\
				Strategic & 30\% & 250 & 75 \\
				\hline
				\textbf{Total} & & & \textbf{175} \\
				\hline
			\end{tabular}
		\end{table}
		
		\vspace{0.1cm}
		\[
		\boxed{\text{VaR}_{95\%} = \$250M}
		\]
		\[
		\boxed{\text{CVaR} = \$275M}
		\]
	\end{frame}
	
	%===========================================================================
	% Q10 - SOLUTION BEHAVIORAL
	%===========================================================================
	\begin{frame}{Q10 Solution: Behavioral Risk}
		\question{Part (d): Behavioral Biases}
		\small
		
		\textbf{Overconfidence}
		Underestimate extreme events
		
		\vspace{0.1cm}
		\textbf{Anchoring}
		Fixate on historical data
		
		\vspace{0.1cm}
		\textbf{Confirmation Bias}
		Seek confirming evidence
		
		\vspace{0.1cm}
		\textbf{Availability Bias}
		Overweight recent events
	\end{frame}
	
	%===========================================================================
	% Q10 - SOLUTION COUNTERWEIGHT
	%===========================================================================
	\begin{frame}{Q10 Solution: ERM Counterweight}
		\question{Part (e): ERM as Counterweight}
		\small
		
		\textbf{1. Independent Oversight}
		Separate risk function, board access
		
		\vspace{0.1cm}
		\textbf{2. Scenario Planning}
		Test tail risk scenarios
		
		\vspace{0.1cm}
		\textbf{3. Diverse Perspectives}
		External experts, dissenting views
		
		\vspace{0.1cm}
		\textbf{4. Decision Protocols}
		Structured approvals, documentation
		
		\vspace{0.1cm}
		\textbf{5. Post-Mortem}
		Review, learn, improve
	\end{frame}
	
	%===========================================================================
	% Q11 - PROBLEM
	%===========================================================================
	\begin{frame}{Q11: Corporate Governance}
		\question{Context}
		\small
		Activist shareholder pressure on governance.
		
		\vspace{0.1cm}
		Current Compensation: Salary \$3M, Bonus \$5M (EPS-based), Equity 200k shares, Options 100k
		
		\vspace{0.1cm}
		\question{Requirements}
		\begin{enumerate}
			\item Explain governance elements.
			\item Analyze agency conflicts.
			\item Evaluate compensation.
			\item Propose reforms.
			\item Discuss activism.
		\end{enumerate}
	\end{frame}
	
	%===========================================================================
	% Q11 - SOLUTION GOVERNANCE
	%===========================================================================
	\begin{frame}{Q11 Solution: Governance Structures}
		\question{Part (a): Governance Structures}
		\small
		
		\textbf{Board}
		Independent majority, independent chair, evaluations, diversity
		
		\vspace{0.1cm}
		\textbf{Committees}
		Audit, Compensation, Nominating, Risk
		
		\vspace{0.1cm}
		\textbf{Compensation}
		Align with long-term value, performance-based, clawbacks
		
		\vspace{0.1cm}
		\textbf{Shareholder Rights}
		Proxy access, majority voting, say-on-pay
	\end{frame}
	
	%===========================================================================
	% Q11 - SOLUTION AGENCY
	%===========================================================================
	\begin{frame}{Q11 Solution: Agency Conflicts}
		\question{Part (b): Agency Conflicts}
		\small
		
		\textbf{1. Shareholder vs Management}
		Compensation, perquisites, risk-taking
		
		\vspace{0.1cm}
		\textbf{2. Controlling vs Minority}
		Related-party transactions, voting rights
		
		\vspace{0.1cm}
		\textbf{3. Debt vs Equity}
		Asset substitution, underinvestment
		
		\vspace{0.1cm}
		\textbf{4. Free Cash Flow (Jensen, 1986)}
		Managers invest in negative NPV projects
	\end{frame}
	
	%===========================================================================
	% Q11 - SOLUTION COMPENSATION
	%===========================================================================
	\begin{frame}{Q11 Solution: Compensation Reform}
		\question{Part (c): Compensation Reform}
		\small
		\begin{table}[h]
			\centering
			\begin{tabular}{|c|c|c|}
				\hline
				\textbf{Component} & \textbf{Current} & \textbf{Proposed} \\
				\hline
				Bonus & \$5M (EPS) & \$4M (Balanced) \\
				Equity & 200k shares & 150k PSUs (ROIC) \\
				\hline
			\end{tabular}
		\end{table}
		
		\vspace{0.1cm}
		\textbf{Balanced Scorecard:}
		Financial, Customer, Internal, Learning
		
		\vspace{0.1cm}
		\textbf{Reforms:}
		Clawbacks, 5-year vesting, no hedging, independent consultant
	\end{frame}
	
	%===========================================================================
	% Q11 - SOLUTION ACTIVISM
	%===========================================================================
	\begin{frame}{Q11 Solution: Shareholder Activism}
		\question{Part (d): Shareholder Activism}
		\small
		
		\textbf{Strategies}
		Board representation, capital return, strategic changes, governance reforms
		
		\vspace{0.1cm}
		\textbf{Impact}
		Increased accountability, transparency, value
		
		\vspace{0.1cm}
		\textbf{Caution}
		Potential short-term focus
		
		\vspace{0.1cm}
		\textbf{Evidence}
		Net positive impact; engage proactively
	\end{frame}
	
	%===========================================================================
	% Q12 - PROBLEM
	%===========================================================================
	\begin{frame}{Q12: ESG Integration}
		\question{Context}
		\small
		Pressure to integrate ESG factors into financial strategy.
		
		\vspace{0.1cm}
		\question{Requirements}
		\begin{enumerate}
			\item Compare shareholder primacy vs stakeholder theory.
			\item Analyze ESG evidence.
			\item Propose ESG framework.
			\item Discuss future of ESG.
		\end{enumerate}
	\end{frame}
	
	%===========================================================================
	% Q12 - SOLUTION THEORY
	%===========================================================================
	\begin{frame}{Q12 Solution: Theory Comparison}
		\question{Part (a): Shareholder vs Stakeholder}
		\small
		
		\textbf{Shareholder Primacy (Friedman, 1970)}
		Only social responsibility is to increase profits
		
		\vspace{0.1cm}
		\textbf{Stakeholder Theory (Freeman, 1984)}
		Create value for all stakeholders (employees, customers, communities, environment)
		
		\vspace{0.1cm}
		\textbf{Enlightened Shareholder Value}
		Shareholder value maximized by considering stakeholders
	\end{frame}
	
	%===========================================================================
	% Q12 - SOLUTION EVIDENCE
	%===========================================================================
	\begin{frame}{Q12 Solution: ESG Evidence}
		\question{Part (b): Empirical Evidence}
		\small
		
		\textbf{Value Creation}
		Positive correlation between ESG and performance
		
		\vspace{0.1cm}
		\textbf{Cost of Capital}
		Lower cost of equity and debt for high ESG firms
		
		\vspace{0.1cm}
		\textbf{Risk Reduction}
		Lower operational, reputational, regulatory risk
		
		\vspace{0.1cm}
		\textbf{Investor Demand}
		Growing demand for ESG investments
	\end{frame}
	
	%===========================================================================
	% Q12 - SOLUTION FRAMEWORK
	%===========================================================================
	\begin{frame}{Q12 Solution: ESG Framework}
		\question{Part (c): ESG Framework}
		\small
		
		\textbf{Environmental}
		Carbon reduction (30\% by 2030), renewable energy, sustainable supply chains
		
		\vspace{0.1cm}
		\textbf{Social}
		Employee welfare, community engagement, diversity
		
		\vspace{0.1cm}
		\textbf{Governance}
		Board oversight, ESG committee, sustainability reporting
	\end{frame}
	
	%===========================================================================
	% Q12 - SOLUTION IMPLEMENTATION
	%===========================================================================
	\begin{frame}{Q12 Solution: Implementation}
		\question{Part (d): Implementation Roadmap}
		\small
		
		\textbf{Year 1:} Assessment, baseline, materiality, stakeholder engagement
		
		\vspace{0.1cm}
		\textbf{Year 2:} Strategy, targets, action plans, pilot programs
		
		\vspace{0.1cm}
		\textbf{Year 3:} Full integration, ESG-linked compensation, reporting
		
		\vspace{0.1cm}
		\textbf{Year 4+:} Continuous improvement, benchmarking, feedback
	\end{frame}
	
	%===========================================================================
	% CONCLUSION
	%===========================================================================
	\begin{frame}{Conclusion: Key Takeaways}
		\question{Summary}
		\small
		
		\textbf{Capital Structure}
		MM foundation, trade-off theory, pecking order
		
		\vspace{0.1cm}
		\textbf{Dividend Policy}
		Irrelevance, bird-in-hand, tax preference, signaling
		
		\vspace{0.1cm}
		\textbf{M\&A}
		DCF valuation, synergies, behavioral biases, integration
		
		\vspace{0.1cm}
		\textbf{Governance}
		ERM framework, agency conflicts, compensation, activism
		
		\vspace{0.1cm}
		\textbf{ESG}
		Stakeholder value, evidence, integration, implementation
	\end{frame}
	
	%===========================================================================
	% FINAL
	%===========================================================================
	\begin{frame}{Final Remarks}
		\begin{center}
			\Large\textcolor{navyblue}{Thank You}
			\vspace{0.2cm}
			
			\small
			\textit{The modern finance professional must be both a quantitative analyst and a behavioral psychologist.}
			\vspace{0.2cm}
			
			\large\textbf{Shivgan Joshi}
		\end{center}
	\end{frame}
	
	%===========================================================================
	% REFERENCES
	%===========================================================================
	\begin{frame}{References}
		\tiny
		\begin{thebibliography}{99}
			\bibitem{beasley2025} Beasley, M. S. (2025). State of risk oversight (16th ed.).
			\bibitem{bis2025} BIS. (2025). Annual economic report 2025.
			\bibitem{fama1970} Fama, E. F. (1970). Efficient capital markets. JF, 25(2).
			\bibitem{fama1993} Fama \& French. (1993). Common risk factors. JFE, 33(1).
			\bibitem{kahneman1979} Kahneman \& Tversky. (1979). Prospect theory. Econometrica, 47(2).
			\bibitem{lo2004} Lo, A. W. (2004). Adaptive markets. JPM, 30(5).
			\bibitem{myers1984} Myers \& Majluf. (1984). Financing decisions. JFE, 13(2).
			\bibitem{modigliani1958} Modigliani \& Miller. (1958). Cost of capital. AER, 48(3).
			\bibitem{sharpe1964} Sharpe, W. F. (1964). CAPM. JF, 19(3).
		\end{thebibliography}
	\end{frame}
	
	%===========================================================================
	% END
	%===========================================================================
\end{document}